
\((\Ch A)_n\) est un \(\lambda\)-idéal.

Soit alors \(x=(n,1+a_1+\ldots+a_n+\ldots)\) un élément d'augmentation \(n\) de \(\Ch A\)~; \(x\) est congru \(\mathrm{mod.}\ (\Ch A)_{n+1}\) à \((n,1+\ldots+a_n)\), et par suite, \(\lambda_{-1}(x)\) est congru à \(\lambda_{-1}(n,1+\ldots+a_n)=(0,1-(n-1)!a_n+\ldots)\), d'où :
\[
  \lambda_{-1}(x)=(0,1-(n-1)!a_n+\ldots).
\]

\textbf{Corollaire 6.6.4.} \uline{Soit} \(\epsilon\) \uline{l'homomorphisme d'augmentation de} \(\Ch A\)~; \uline{alors}, \uline{pour tout} \(x\in\Ch A\), \uline{on a} :
\[
  \gamma^n(x-\epsilon(x))
  =
  (0,1+(-1)^{n-1}(n-1)!a_n+\ldots)
\]

Par définition, \(\gamma^n(x-\epsilon(x))=(-1)^n\sum_{i=0}^n(-1)^i\lambda(x+n-\epsilon(x))\). Comme \(x+n-\epsilon(x)\) est d'augmentation \(n\) on peut lui appliquer le corollaire 6.6.2, d'où le résultat.

\textbf{Corollaire 6.6.5.} \uline{Avec les notations de} 3.10, \(\Fil_n\Ch A\subset(\Ch A)_n\).

D'après le corollaire précédent, \(\gamma^n(x-\epsilon(x))\in(\Ch A)_n\). Le corollaire résulte alors de ce que les deux filtrations sont des filtrations d'anneau, et de ce que les \(\gamma^n(x-\epsilon(x))\), pour \(n\) et \(x\) variables, engendrent la filtration de \(\lambda\)-anneau.

\textbf{Corollaire 6.6.7.} \uline{Soit} \(A\) \uline{une} \(K\)-\uline{algèbre graduée de caractéristique nulle}, \uline{à degrés bornés}~; \uline{alors les deux filtrations définies sur} \(\Ch A\) \uline{sont} \uline{identiques}.

Supposons que \(\Fil^pA=0\), et soit \(x=(0,1+a_n+\ldots+a_p)\) un élément de \((\Ch A)_n\). Posons \(x'=\left(0,1+\dfrac{(-1)^{n-1}}{(n-1)!}a_n+a_{n+1}+\ldots+a_p\right)\)~; on a alors :
\(\gamma^n(x')=(0,1+a_n+\ldots+a'_p)\), car \(x'\) est d'augmentation nulle. Par suite, 
